\subsubsection{Moduli Stabilisation for ADD and the Dark Dimension Scenario --- arXiv:2606.19440}

\textbf{Authors:} Andreas P. Braun, Michele Cicoli, Riccardo Milioli, Roberto Valandro

\paragraph{主な主張}
Type IIB Large Volume Scenario の K3-fibred Calabi--Yau において、string loop と高階微分補正で K3 fibre を比較的小さく安定化し、complex structure の Tyurin degeneration で $\mathbb{P}^1$ base を異方化することで、2つの大余剰次元を持つ ADD と1つの大余剰次元を持つ Dark Dimension を同一の枠組みで実現できることを示した \cite{Braun:2026ADD}。

\paragraph{新規性}
ADD と Dark Dimension を別々の compactification としてではなく、K\"ahler moduli と complex structure moduli の異方的安定化を組み合わせた統一的な機構として記述し、brane setup・tadpole cancellation・moduli stabilisation を備えた具体的な Calabi--Yau orientifold 例まで構成した。

\paragraph{位置付け}
Large extra dimensions の string embedding として、compactification の volume だけでなく shape moduli が実効的な大余剰次元の数を変え得ることを具体化しており、余剰次元の「大きさ」に加えて「形状」の観測的効果を考える際の理論的基盤となる。
